% Source: physical PDF p. 281 (printed p. 258). Coverage: Figure 20.1 and complete caption. Uncertainty: none; geometry, momentum labels, arrows, hatching, divider, equality sign, plus sign, and disk labels reconstructed from the rendered source.
\begin{figure}[htbp]
 \centering
 \begin{tikzpicture}[x=1cm,y=1cm,line cap=round,line join=round,
   every node/.style={inner sep=1pt}]
  \tikzset{
   flow/.style={line width=0.75pt,-{Stealth[length=1.7mm,width=1.25mm]}},
   rail/.style={line width=0.75pt}
  }

  % A cross-hatched amplitude disk.
  \newcommand{\GammaDisk}[2]{%
   \fill[white] (#1,#2) circle[radius=0.72];
   \begin{scope}
    \clip (#1,#2) circle[radius=0.72];
    \foreach \d in {-1.35,-1.05,...,1.35}
     \draw[line width=0.45pt]
      ({#1-0.95},{#2+\d-0.42}) -- ({#1+0.95},{#2+\d+0.42});
   \end{scope}
   \draw[line width=0.9pt] (#1,#2) circle[radius=0.72];
  }
  \newcommand{\IDisk}[2]{%
   \GammaDisk{#1}{#2}
   \draw[line width=0.8pt] (#1,{#2-0.72}) -- (#1,{#2+0.72});
  }

  % Left-hand amplitude Gamma.
  \draw[rail] (0.00,0.55) -- (3.15,0.55);
  \draw[rail] (0.00,-0.55) -- (3.15,-0.55);
  \draw[flow] (0.15,0.55) -- (0.92,0.55);
  \draw[flow] (0.15,-0.55) -- (0.92,-0.55);
  \draw[flow] (2.26,0.55) -- (3.00,0.55);
  \draw[flow] (2.26,-0.55) -- (3.00,-0.55);
  \GammaDisk{1.58}{0}
  \node[above] at (0.42,0.66) {$k$};
  \node[above] at (2.76,0.66) {$k'$};
  \node[below] at (0.50,-0.66) {$p-k$};
  \node[below] at (2.66,-0.66) {$p-k'$};
  \node at (1.58,-1.13) {$\Gamma$};

  \node[font=\Large] at (3.68,0) {$=$};

  % Irreducible kernel I.
  \draw[rail] (4.18,0.55) -- (7.25,0.55);
  \draw[rail] (4.18,-0.55) -- (7.25,-0.55);
  \draw[flow] (4.32,0.55) -- (5.08,0.55);
  \draw[flow] (4.32,-0.55) -- (5.08,-0.55);
  \draw[flow] (6.38,0.55) -- (7.12,0.55);
  \draw[flow] (6.38,-0.55) -- (7.12,-0.55);
  \IDisk{5.72}{0}
  \node[above] at (4.59,0.66) {$k$};
  \node[above] at (6.88,0.66) {$k'$};
  \node[below] at (4.67,-0.66) {$p-k$};
  \node[below] at (6.79,-0.66) {$p-k'$};
  \node at (5.72,-1.13) {$I$};

  \node[font=\Large] at (7.70,0) {$+$};

  % Iterated kernel I joined to Gamma by the k'' bridge.
  \draw[rail] (8.14,0.55) -- (14.65,0.55);
  \draw[rail] (8.14,-0.55) -- (14.65,-0.55);
  \draw[flow] (8.26,0.55) -- (9.02,0.55);
  \draw[flow] (8.26,-0.55) -- (9.02,-0.55);
  \draw[flow] (10.34,0.55) -- (11.24,0.55);
  \draw[flow] (10.34,-0.55) -- (11.24,-0.55);
  \draw[flow] (12.58,0.55) -- (14.50,0.55);
  \draw[flow] (12.58,-0.55) -- (14.50,-0.55);
  \IDisk{9.66}{0}
  \GammaDisk{11.92}{0}
  \node[above] at (8.55,0.66) {$k$};
  \node[above] at (10.78,0.66) {$k''$};
  \node[above] at (14.20,0.66) {$k'$};
  \node[below] at (8.62,-0.66) {$p-k$};
  \node[below] at (10.78,-0.66) {$p-k''$};
  \node[below] at (14.08,-0.66) {$p-k'$};
  \node at (9.66,-1.13) {$I$};
  \node at (11.92,-1.13) {$\Gamma$};
 \end{tikzpicture}
 \renewcommand{\thefigure}{20.1}
 \caption{Diagrammatic representation of the integral equation (20.2.3). The cross-hatched disks marked $\Gamma$ represent the sum of all connected Feynman diagrams with the indicated external lines, while the cross-hatched disks marked $I$, which are divided by a vertical line, represent the sum of all connected diagrams for which the lines on the left cannot be separated from those on the right by cutting through any pair of external lines.}
 \label{fig:20.1}
\end{figure}
